\subsection{Post-hoc analyses requested in review}\label{subsec:posthoc}

{\color{blue}The analyses in this subsection were requested by the reviewers and specified after the registered evaluation, in two dated addenda to the protocol that were committed to the code repository before they were run. They are descriptive, are not part of the registered tests, and are reported in full; none of them changes the registered forecaster, agents, selection, or tests.

\paragraph{Relation layers and unrestricted inter-stock attention}
Table~\ref{tab:xablation} reports MGAT variants that remove one relation layer from the edge set of Eq.~\eqref{eq:Afused_re} or admit all ordered stock pairs, trained in the same four folds with three seeds. Because the positive entries of the Kendall graph alone cover 90\% (banking) and 82\% (healthcare) of the ordered pairs, and every shareholder edge is also a Kendall edge, removing the shareholder or Q\&A layer changes few or no admissible pairs: the variant without the shareholder layer has exactly the neighbourhoods of MGAT in both sectors, and so does the variant without the Q\&A layer in healthcare. These retrainings use the same code, data, and seeds as the registered runs but a different number of CPU threads, and therefore a different order of floating-point operations; they measure the run-to-run variation of the training itself: their IC$_2$ differs from that of the registered MGAT by up to 0.048 in a development year (banking, 2023) and by 0.006 or less in 2025. Removing the Kendall layer restricts attention to firms linked by shareholding or co-mentions (75\% and 36\% of the pairs); this variant has the highest 2025 IC$_2$ of the variants in both sectors (0.021 and 0.064) but not in the development years, and unrestricted attention over all pairs has a lower IC$_2$ than MGAT in banking and a slightly higher one in healthcare. In 2025, all paired daily IC$_2$ differences between MGAT and these variants lie between $-0.012$ and 0.011, with 95\% block-bootstrap intervals that include zero, and the high and low forecasts differ by less than one percentage point. The relation layers therefore mainly exclude a minority of stock pairs, and neither removing a layer nor admitting all pairs changes the forecasts beyond the variation between training runs. Unrestricted attention has the same number of parameters as MGAT and adds only the coefficients of the excluded pairs; its mean training time per run in fold 2025 (61--77~s, under a heavier machine load than the registered runs) is comparable with that of MGAT (62--65~s), so its computational overhead is negligible for fifteen stocks.

\paragraph{Graph-attention coefficients}
Table~\ref{tab:attention} summarises the attention coefficients of the registered MGAT on the 2025 inputs. In banking, the coefficients are close to uniform over the admissible neighbours in all three layers (entropy 94--100\% of that of uniform attention), and pairs linked by shareholding or co-mentions receive the same attention as other pairs (ratios 0.97--1.00). In healthcare, the first layer gives pairs co-mentioned in investor questions 24\% and pairs with common shareholders 10\% more weight than the average admissible neighbour and pairs linked only by trading-value correlation 6\% less, whereas the second and third layers are nearly uniform. The first-layer coefficients vary from day to day (median coefficient of variation 0.25 and 0.33) and agree only moderately across seeds (correlation 0.52 and 0.41), whereas those of the deeper layers are stable. The learned attention therefore reflects the qualitative relations weakly, and only in the first layer and in the sector where the forecasts carry ranking information; it does not single out a small set of economically interpretable links.

\paragraph{Depth of the graph and Transformer components}
Table~\ref{tab:xablation} also reports variants with one graph-attention layer (64 units; 42{,}526 parameters), two layers (32 and 64 units; 44{,}414 parameters), and two Transformer encoder blocks (84{,}790 parameters). In 2025, their IC$_2$ is 0.002--0.015 in banking and 0.044--0.065 in healthcare, against 0.014 and 0.052 for MGAT; the paired daily differences from MGAT lie between $-0.013$ and 0.012, with 95\% intervals that include zero, and their high and low forecasts differ from those of MGAT by at most 2.3 percentage points. Over the development years, the three-layer MGAT has the highest IC$_2$ in banking (0.057, against 0.002--0.041), and the two-layer variant the highest in healthcare (0.026, against 0.018 for MGAT). The depth of the graph encoder therefore does not change the conclusions, and a second Transformer block, which doubles the number of parameters, does not improve the forecasts.

\paragraph{Direction of the forecasts}
As a further decision-relevant measure, we computed the share of stock-days on which the sign of $\hat q_{t,i}$, or of $\hat q_{t,i}$ relative to the day's cross-sectional mean, matches that of the realised $q_{t,i}$. In 2025, MGAT has hit rates of 50.8\% and 50.2\% in banking and 50.5\% and 51.4\% in healthcare. Across all forecasters of Tables~\ref{tab:forecast_eval} and~\ref{tab:xablation}, the hit rates lie between 48.4\% and 52.6\% and the relative hit rates between 48.9\% and 52.1\%; the largest deviation from 50\% has $t=2.0$ (one-layer variant, healthcare, relative sign), and the other 51 hit rates have $|t|\le1.9$. Because the cross-sectional rank correlation uses the full ordering rather than signs, it is the more informative measure of the weak signal.

\paragraph{Allocation baselines}
Panel~A of Table~\ref{tab:posthoc_port} adds long-only allocation rules that use only the 252-day covariance matrix of the RL state, with the same ledger and costs. In banking, the minimum-variance portfolio attains a Sharpe ratio of 1.27 (monthly rebalancing) and 1.25 (daily) in 2025 under F1, with an annualised return of 19.9\% and 19.3\%, compared with 0.61 for daily equal weighting and 0.54--0.58 for the RL agents of every forecast source; it also exceeded daily equal weighting and the mean of the selected policies in the selection year 2024 (1.85, against 1.70 and 1.64--1.69). In healthcare, minimum variance (1.19) and inverse volatility (1.20) are close to equal weighting (1.23) and to the MGAT agents (1.20). A rule that requires no forecast thus outperformed every learned policy in banking in 2025. Because the agents observe the same covariance matrix, this result shows that they also did not exploit the risk information in their state; their allocations remain close to equal weighting (Table~\ref{tab:policy}).

\paragraph{Volatility regimes within 2025}
Panel~B splits the 2025 trading days at the median of the trailing 20-day volatility of daily equal weighting (121 days each). The paired difference between the MGAT and naive agents is $-0.008$ $[-0.115, 0.087]$ on high-volatility and $-0.068$ $[-0.144, 0.004]$ on low-volatility days in banking, and $-0.011$ $[-0.135, 0.133]$ and 0.107 $[-0.065, 0.267]$ in healthcare. No interval excludes zero; the positive healthcare difference of the full year comes from the calmer half of the year.

\paragraph{PPO learning rate and discount factor}
Panel~C of Table~\ref{tab:posthoc_port} retrains the PPO agents of the MGAT and naive sources (seeds 1--5) with one hyperparameter changed and otherwise the registered training and selection. Across the five settings, the mean Sharpe ratio of the MGAT agents lies between 0.48 and 0.57 in banking and between 1.16 and 1.21 in healthcare, against 0.46--0.57 and 1.14--1.20 for the naive agents. Seven of the eight changed settings give paired differences whose 90\% intervals include zero. With $\gamma=0.999$ in banking, the MGAT agents have a higher mean Sharpe ratio (0.066, $[0.020, 0.109]$, one-sided $p=0.008$), but the MGAT agent is higher for only two of the five seeds, and the difference is driven by large gaps in a minority of pairs (the standard deviation of the naive agents' Sharpe ratios is 0.13, against 0.03 for the MGAT agents); it is not significant after a Bonferroni correction for the eight comparisons (threshold 0.006). The conclusions of the registered PPO comparison therefore do not depend on the learning rate or the discount factor within these ranges.

\paragraph{Market structure}
Daily returns co-move strongly in banking: the mean pairwise correlation is 0.56--0.63 in 2022--2025, and the first principal component explains 60--65\% of the standardised variance. In healthcare, the correlation is 0.30--0.51, the first component explains 35--55\%, and the mean cross-sectional standard deviation of daily returns is 1.8--2.4\%, compared with 0.7--0.8\% in banking. The larger cross-sectional dispersion leaves more room for ranking stocks, which is consistent with the significant 2025 IC$_2$ in healthcare and its absence in banking; with two sectors, this association is descriptive.\par}

\input{tab_xablation}
\input{tab_attention}
\input{tab_posthoc_port}
